\input cwebmac

\N{1}{1}Hello X11.

This is a small graphical program written in standard~C using
Donald Knuth's literate programming paradigm.  It opens a single
top-level window via the X~Window~System and draws two clickable
buttons inside it:

\medskip
$${\vbox{
\item{} {\bf Hello} --- prints
\.{Hello World, and the current local time.}
(followed by the current local time) to standard output.

\item{} {\bf Exit}  --- closes the window and terminates the program.
}}$$

\medskip\noindent
The window may also be closed via the standard window-manager close
gesture (the \.{WM\_DELETE\_WINDOW} protocol), in which case the program
exits cleanly as if {\bf Exit} had been clicked.  All graphics and event
handling are performed with raw \.{Xlib} calls (\.{libX11}); no toolkit
(Xt, Motif, GTK, Qt) is involved.

\fi

\M{2}{\bf How it works.}
The program is a textbook Xlib client.  It performs four phases in
sequence:

\medskip\item{1.}  Connect to the X~display server and discover screen
defaults (root window, default colours, default font).

\item{2.}  Create a top-level window, install a graphics context, and
request the events of interest (exposure, button presses, structure
notifications).

\item{3.}  Enter an event loop: each event is dispatched to a handler
that either repaints the buttons (\.{Expose}), checks for a button
click (\.{ButtonPress}), or terminates the loop
(\.{ClientMessage} with \.{WM\_DELETE\_WINDOW}).

\item{4.}  On exit, release the graphics context and close the display.
\medskip

\fi

\M{3}{\bf Literate programming.}
Donald Knuth introduced {\it literate programming\/} in 1984 as a way of
writing software that is meant to be read by human beings first and
executed by computers second.  Rather than annotating code with
comments, a literate program interweaves prose and code in a single
source document.  The prose explains the {\it why\/}---the motivation,
the design decisions, the mathematical reasoning---while the code
expresses the {\it how}.  The two live together in one file
(conventionally given the extension \.{.w}) and are separated only at
build time by two companion tools: \.{ctangle} and \.{cweave}.

\fi

\M{4}{\bf ctangle.}
\.{ctangle} is the {\it tangling\/} tool.  It reads a \.{.w} source
file and extracts the C~code sections, assembling them in the order
dictated by named chunk references rather than the order in which they
appear in the document.  The result is a plain \.{.c} file that a
standard C~compiler can process without any knowledge of literate
programming.  In this project, running
$$\.{ctangle hello-x11.w}$$
produces \.{hello-x11.c}, which is then compiled with \.{cc} and linked
against \.{libX11} to create the \.{hello-x11} executable.  The
generated \.{.c} file should be treated as a build artefact: the
\.{.w} file is the true source of record.

\fi

\M{5}{\bf cweave.}
\.{cweave} is the {\it weaving\/} tool.  It reads the same \.{.w} source
and produces a \.{.tex} file formatted for \.{pdftex} using the
\.{cwebmac} macro package.  \.{cweave} pretty-prints all C~code with
bold keywords, italic identifiers, and cross-references, and numbers
every named chunk so the reader can follow the program's logical
structure independently of its physical layout.  An index of identifiers
and a table of contents are generated automatically.  Running
$$\.{cweave hello-x11.w}$$
produces \.{hello-x11.tex}; running \.{pdftex} on that file yields the
typeset documentation you are reading now.

\fi

\M{6}{\bf Compilation.}  After tangling with \.{ctangle}:
$$\.{cc -O2 -o hello-x11 hello-x11.c -lX11}$$
On systems where \.{X11/Xlib.h} is not on the default include path, add
\.{-I/usr/X11R6/include -L/usr/X11R6/lib} (Linux/BSD) or
\.{-I/opt/X11/include -L/opt/X11/lib} (macOS with XQuartz).

\fi

\M{7}{\bf Program structure.}  The top-level arrangement of the tangled
output is:

\Y\B\X9:Includes\X\6
\X10:Constants\X\6
\X13:Type definitions\X\6
\X14:Globals\X\6
\X17:Time helper\X\6
\X19:Hit testing\X\6
\X21:Drawing routines\X\6
\X25:Event handlers\X\6
\X29:Display setup\X\6
\X35:Display teardown\X\6
\X37:Main function\X\par
\fi

\N{1}{8}Includes and constants.

\fi

\M{9}Three standard headers cover the C~runtime; \.{Xlib.h} is the only
X~Window header needed for this program (the higher-level \.{Xutil.h}
is pulled in for the \.{XSizeHints} structure used to suggest a fixed
window size to the window manager).

\Y\B\4\X9:Includes\X${}\E{}$\6
\8\#\&{include} \.{<stdio.h>}\6
\8\#\&{include} \.{<stdlib.h>}\6
\8\#\&{include} \.{<string.h>}\6
\8\#\&{include} \.{<time.h>}\6
\8\#\&{include} \.{<X11/Xlib.h>}\6
\8\#\&{include} \.{<X11/Xutil.h>}\par
\U7.\fi

\M{10}The window dimensions and the geometry of the two buttons are fixed
at compile time --- the program is not designed to be resized.  The
buttons are stacked vertically with a uniform margin.

\Y\B\4\X10:Constants\X${}\E{}$\6
\8\#\&{define} \.{WIN\_W}\5\T{320}\C{ top-level window width in pixels
}\6
\8\#\&{define} \.{WIN\_H}\5\T{160}\C{ top-level window height in pixels
}\6
\8\#\&{define} \.{BTN\_W}\5\T{120}\C{ button width in pixels
}\6
\8\#\&{define} \.{BTN\_H}\5\T{40}\C{ button height in pixels                  }%
\6
\8\#\&{define} \.{MARGIN}\5\T{20}\C{ margin around and between buttons        }%
\6
\8\#\&{define} \.{WIN\_TITLE}\5\.{"Hello\ X11"}\par
\A11.
\U7.\fi

\M{11}Each button has its origin at the top-left corner of its rectangle.
\PB{\.{HELLO\_X}}/\PB{\.{HELLO\_Y}} place the {\bf Hello} button in the upper
region
of the window; \PB{\.{EXIT\_X}}/\PB{\.{EXIT\_Y}} place the {\bf Exit} button
directly
below it.  Both buttons are horizontally centred.

\Y\B\4\X10:Constants\X${}\mathrel+\E{}$\6
\8\#\&{define} \.{HELLO\_X}\5${}((\.{WIN\_W}-\.{BTN\_W})/\T{2}){}$\6
\8\#\&{define} \.{HELLO\_Y}\5\.{MARGIN}\6
\8\#\&{define} \.{EXIT\_X}\5${}((\.{WIN\_W}-\.{BTN\_W})/\T{2}){}$\6
\8\#\&{define} \.{EXIT\_Y}\5${}(\.{MARGIN}+\.{BTN\_H}+\.{MARGIN}){}$\par
\fi

\N{1}{12}Data structures.

\fi

\M{13}A \PB{\\{Button}} record bundles the four geometry fields and the label
string for a single clickable rectangle.  Storing the geometry in a
struct (rather than chasing six bare \.{\#define}s through every event
handler) keeps the hit-test routine to a single line.

\Y\B\4\X13:Type definitions\X${}\E{}$\6
\&{typedef} \&{struct} ${}\{{}$\1\6
\&{int} \|x${},\39\|y{}$;\C{ top-left corner, in window coordinates    }\6
\&{int} \|w${},\39\|h{}$;\C{ width and height in pixels                }\6
\&{const} \&{char} ${}{*}\\{label}{}$;\C{ null-terminated text drawn inside
    }\2\6
${}\}{}$ \&{Button};\par
\U7.\fi

\M{14}Two globals hold the X~Window state: the open display connection and
the graphics context used for every draw call.  They are initialised
by \PB{\\{display\_open}} and released by \PB{\\{display\_close}}; in between
they are
read by the drawing and event-handling routines.  Keeping them at file
scope avoids threading them through every callback as parameters.

\Y\B\4\X14:Globals\X${}\E{}$\6
\&{static} \\{Display}${}{*}\\{g\_dpy}\K\NULL{}$;\C{ connection to the X~server
        }\6
\&{static} \.{GC}\\{g\_gc}${}\K\NULL{}$;\C{ graphics context for all drawing
}\6
\&{static} \\{Window}\\{g\_win}${}\K\T{0}{}$;\C{ the top-level window
    }\6
\&{static} \\{Atom}\\{g\_wm\_delete}${}\K\T{0}{}$;\C{ \.{WM\_DELETE\_WINDOW}
client atom }\par
\A15.
\U7.\fi

\M{15}The two buttons are file-scope so that \PB{\\{draw\_buttons}} and \PB{%
\\{hit\_test}}
can iterate over them without rebuilding the array on every event.
The order in this array also determines paint order (top to bottom).

\Y\B\4\X14:Globals\X${}\mathrel+\E{}$\6
\&{static} \&{const} \&{Button} \\{g\_buttons}[\T{2}]${}\K\{\{\.{HELLO\_X},\39%
\.{HELLO\_Y},\39\.{BTN\_W},\39\.{BTN\_H},\39\.{"Hello"}\},\39\{\.{EXIT\_X},\39%
\.{EXIT\_Y},\39\.{BTN\_W},\39\.{BTN\_H},\39\.{"Exit"}\}\};{}$\6
\8\#\&{define} \.{N\_BUTTONS}\5((\&{int})(\&{sizeof} \\{g\_buttons}${}/{}$%
\&{sizeof} \\{g\_buttons}[\T{0}]))\par
\fi

\N{1}{16}Time helper.

\fi

\M{17}\PB{\\{format\_local\_time}} writes the current local time into \PB{%
\\{out}} as
\.{YYYY-MM-DD HH:MM:SS}.  \PB{\\{out}} must hold at least 20 bytes
(19~characters plus a null terminator).  The format string is
deliberately unambiguous --- ISO~8601 date plus 24-hour time --- so
the line printed by the {\bf Hello} button is locale-independent.

\Y\B\4\X17:Time helper\X${}\E{}$\6
\1\1\&{static} \&{void} \\{format\_local\_time}(\&{char} ${}{*}\\{out},\39{}$%
\&{size\_t} \\{max\_len})\2\2\6
${}\{{}$\1\6
\&{time\_t} \\{now}${}\K\\{time}(\NULL);{}$\6
\&{struct} \\{tm} ${}{*}\\{lt}\K\\{localtime}({\AND}\\{now});{}$\7
\&{if} (\\{lt})\1\5
${}\\{strftime}(\\{out},\39\\{max\_len},\39\.{"\%Y-\%m-\%d\ \%H:\%M:\%S"},\39%
\\{lt});{}$\2\6
\&{else}\1\5
${}\\{snprintf}(\\{out},\39\\{max\_len},\39\.{"(time\ unavailable)"});{}$\2\6
\4${}\}{}$\2\par
\U7.\fi

\N{1}{18}Hit testing.

\fi

\M{19}\PB{\\{hit\_test}} returns the index of the button containing the point
$(x,y)$, or $-1$ if no button is hit.  The buttons do not overlap, so
the first match wins.

\Y\B\4\X19:Hit testing\X${}\E{}$\6
\1\1\&{static} \&{int} \\{hit\_test}(\&{int} \|x${},\39{}$\&{int} \|y)\2\2\6
${}\{{}$\1\6
\&{for} (\&{int} \|i${}\K\T{0};{}$ ${}\|i<\.{N\_BUTTONS};{}$ ${}\|i\PP){}$\5
${}\{{}$\1\6
\&{const} \&{Button} ${}{*}\|b\K{\AND}\\{g\_buttons}[\|i];{}$\7
\&{if} ${}(\|x\G\|b\MG\|x\W\|x<\|b\MG\|x+\|b\MG\|w\W\|y\G\|b\MG\|y\W\|y<\|b\MG%
\|y+\|b\MG\|h){}$\1\5
\&{return} \|i;\2\6
\4${}\}{}$\2\6
\&{return} ${}{-}\T{1};{}$\6
\4${}\}{}$\2\par
\U7.\fi

\N{1}{20}Drawing routines.

\fi

\M{21}\PB{\\{draw\_button}} paints a single button: an outlined rectangle with
the
label centred horizontally and vertically.  The label is centred by
querying the font metrics of the current graphics context with
\PB{\\{XQueryTextExtents}}; this avoids hard-coding a font width and keeps the
labels centred even if the user installs a different default font.

\Y\B\4\X21:Drawing routines\X${}\E{}$\6
\1\1\&{static} \&{void} \\{draw\_button}(\&{const} \&{Button} ${}{*}\|b)\2\2{}$%
\6
${}\{{}$\1\6
${}\\{XDrawRectangle}(\\{g\_dpy},\39\\{g\_win},\39\\{g\_gc},\39\|b\MG\|x,\39\|b%
\MG\|y,\39\|b\MG\|w,\39\|b\MG\|h);{}$\6
\X22:Centre and draw the button label\X\6
\4${}\}{}$\2\par
\A23.
\U7.\fi

\M{22}The text-extents query reports the label's bounding box in the
current font.  The horizontal centre is straightforward; the vertical
centre uses the font's ascent so the label sits on a baseline that
splits the button cleanly.

\Y\B\4\X22:Centre and draw the button label\X${}\E{}$\6
\&{int} \\{dir}${},\39\\{ascent},\39\\{descent};{}$\7
\\{XCharStruct}\\{ext};\7
\&{int} \\{len}${}\K(\&{int})\,\\{strlen}(\|b\MG\\{label});{}$\7
${}\\{XQueryTextExtents}(\\{g\_dpy},\39\\{XGContextFromGC}(\\{g\_gc}),\39\|b\MG%
\\{label},\39\\{len},\39{\AND}\\{dir},\39{\AND}\\{ascent},\39{\AND}\\{descent},%
\39{\AND}\\{ext});{}$\7
\&{int} \\{tx}${}\K\|b\MG\|x+(\|b\MG\|w-\\{ext}.\\{width})/\T{2};{}$\6
\&{int} \\{ty}${}\K\|b\MG\|y+(\|b\MG\|h+\\{ascent}-\\{descent})/\T{2};{}$\7
${}\\{XDrawString}(\\{g\_dpy},\39\\{g\_win},\39\\{g\_gc},\39\\{tx},\39\\{ty},%
\39\|b\MG\\{label},\39\\{len}){}$;\par
\U21.\fi

\M{23}\PB{\\{draw\_buttons}} repaints every button.  It is called on every
\.{Expose} event so the window can recover from being uncovered or
resized by the window manager.

\Y\B\4\X21:Drawing routines\X${}\mathrel+\E{}$\6
\1\1\&{static} \&{void} \\{draw\_buttons}(\&{void})\2\2\6
${}\{{}$\1\6
\&{for} (\&{int} \|i${}\K\T{0};{}$ ${}\|i<\.{N\_BUTTONS};{}$ ${}\|i\PP){}$\1\5
${}\\{draw\_button}({\AND}\\{g\_buttons}[\|i]);{}$\2\6
\4${}\}{}$\2\par
\fi

\N{1}{24}Event handlers.

\fi

\M{25}\PB{\\{on\_hello\_click}} prints the greeting and the current local time
to
standard output.  \PB{\\{fflush}} is called explicitly because \.{stdout} is
typically line-buffered when attached to a terminal but
{\it block\/}-buffered when redirected to a pipe or file --- without
the flush, a user redirecting the output to \.{tee} would see nothing
until the program exited.

\Y\B\4\X25:Event handlers\X${}\E{}$\6
\1\1\&{static} \&{void} \\{on\_hello\_click}(\&{void})\2\2\6
${}\{{}$\1\6
\&{char} \\{ts}[\T{32}];\7
${}\\{format\_local\_time}(\\{ts},\39{}$\&{sizeof} \\{ts});\6
${}\\{printf}(\.{"Hello\ World,\ and\ th}\)\.{e\ current\ local\ time}\)\.{:\ %
\%s\\n"},\39\\{ts});{}$\6
\\{fflush}(\\{stdout});\6
\4${}\}{}$\2\par
\As26\ET27.
\U7.\fi

\M{26}\PB{\\{on\_button\_press}} dispatches a \.{ButtonPress} event to the
appropriate handler.  Returning non-zero asks the main loop to exit;
this is how the {\bf Exit} button terminates the program.  Clicks
outside any button are silently ignored.

\Y\B\4\X25:Event handlers\X${}\mathrel+\E{}$\6
\1\1\&{static} \&{int} \\{on\_button\_press}(\&{const} \\{XButtonEvent}${}{*}%
\\{ev})\2\2{}$\6
${}\{{}$\1\6
\&{int} \\{idx}${}\K\\{hit\_test}(\\{ev}\MG\|x,\39\\{ev}\MG\|y);{}$\7
\&{if} ${}(\\{idx}<\T{0}){}$\1\5
\&{return} \T{0};\2\6
\&{if} ${}(\\{strcmp}(\\{g\_buttons}[\\{idx}].\\{label},\39\.{"Hello"})\E%
\T{0}){}$\5
${}\{{}$\1\6
\\{on\_hello\_click}(\,);\6
\&{return} \T{0};\6
\4${}\}{}$\2\6
\&{if} ${}(\\{strcmp}(\\{g\_buttons}[\\{idx}].\\{label},\39\.{"Exit"})\E%
\T{0}){}$\5
${}\{{}$\1\6
\&{return} \T{1};\6
\4${}\}{}$\2\6
\&{return} \T{0};\6
\4${}\}{}$\2\par
\fi

\M{27}\PB{\\{on\_client\_message}} handles the \.{WM\_DELETE\_WINDOW} protocol:
when the user clicks the window manager's close decoration, the server
delivers a \.{ClientMessage} carrying the \PB{\\{g\_wm\_delete}} atom.  We
return non-zero to break out of the event loop so the close button
behaves identically to the in-window {\bf Exit} button.

\Y\B\4\X25:Event handlers\X${}\mathrel+\E{}$\6
\1\1\&{static} \&{int} \\{on\_client\_message}(\&{const} %
\\{XClientMessageEvent}${}{*}\\{ev})\2\2{}$\6
${}\{{}$\1\6
\&{if} ${}((\\{Atom})\\{ev}\MG\\{data}.\|l[\T{0}]\E\\{g\_wm\_delete}){}$\1\5
\&{return} \T{1};\2\6
\&{return} \T{0};\6
\4${}\}{}$\2\par
\fi

\N{1}{28}Display setup.

\fi

\M{29}\PB{\\{display\_open}} performs all four pre-loop X~Window steps: open
the
display, create the window, create the graphics context, and register
for the \.{WM\_DELETE\_WINDOW} protocol.  The body is split into four
sub-modules so no single chunk exceeds twenty-four lines.  The function
returns 1 on success or 0 on any failure, with a diagnostic on
\.{stderr}.

\Y\B\4\X29:Display setup\X${}\E{}$\6
\1\1\&{static} \&{int} \\{display\_open}(\&{void})\2\2\6
${}\{{}$\1\6
${}\X30:Open the X display\X\X31:Create the top-level window\X\X32:Create the
graphics context\X\X33:Register for window-manager close\X\\{XMapWindow}(\\{g%
\_dpy},\39\\{g\_win});{}$\6
\&{return} \T{1};\6
\4${}\}{}$\2\par
\U7.\fi

\M{30}The display name is read from the \PB{\.{DISPLAY}} environment variable
when
the argument to \PB{\\{XOpenDisplay}} is \PB{$\NULL$}.  A failure here
typically
means \.{DISPLAY} is unset or the user lacks permission to connect to
the named server.

\Y\B\4\X30:Open the X display\X${}\E{}$\6
$\\{g\_dpy}\K\\{XOpenDisplay}(\NULL);{}$\6
\&{if} ${}(\R\\{g\_dpy}){}$\5
${}\{{}$\1\6
${}\\{fputs}(\.{"Error:\ cannot\ open\ }\)\.{X\ display\ "}\.{"(is\ \$DISPLAY\
set?)\\}\)\.{n"},\39\\{stderr});{}$\6
\&{return} \T{0};\6
\4${}\}{}$\2\par
\U29.\fi

\M{31}The window is created as a child of the screen's root window, with a
white background and a one-pixel black border.  We immediately call
\PB{\\{XSelectInput}} to subscribe to the three event masks the program needs:
\.{ExposureMask} (initial paint and repaints),
\.{ButtonPressMask} (mouse clicks), and
\.{StructureNotifyMask} (the window-manager close gesture is delivered
as a structure-notify-class \.{ClientMessage}).
The title is set with \PB{\\{XStoreName}} so window managers display
\.{Hello X11} in the title bar.

\Y\B\4\X31:Create the top-level window\X${}\E{}$\6
\&{int} \\{scr}${}\K\\{DefaultScreen}(\\{g\_dpy});{}$\7
${}\\{Window}\\{root}\K\\{RootWindow}(\\{g\_dpy},\39\\{scr});{}$\6
${}\\{g\_win}\K\\{XCreateSimpleWindow}(\\{g\_dpy},\39\\{root},\39\T{0},\39%
\T{0},\39\.{WIN\_W},\39\.{WIN\_H},\39\T{1},\39\\{BlackPixel}(\\{g\_dpy},\39%
\\{scr}),\39\\{WhitePixel}(\\{g\_dpy},\39\\{scr}));{}$\6
${}\\{XStoreName}(\\{g\_dpy},\39\\{g\_win},\39\.{WIN\_TITLE});{}$\6
${}\\{XSelectInput}(\\{g\_dpy},\39\\{g\_win},\39\\{ExposureMask}\OR%
\\{ButtonPressMask}\OR\\{StructureNotifyMask}){}$;\par
\U29.\fi

\M{32}The graphics context is created against the window's drawable.  We
pre-set the foreground colour to black --- the default is
implementation-defined.  All draw calls (rectangles and strings) reuse
this single GC; there is no need for separate contexts since the
program draws in only one colour and one font.

\Y\B\4\X32:Create the graphics context\X${}\E{}$\6
$\\{g\_gc}\K\\{XCreateGC}(\\{g\_dpy},\39\\{g\_win},\39\T{0},\39\NULL);{}$\6
\&{if} ${}(\R\\{g\_gc}){}$\5
${}\{{}$\1\6
${}\\{fputs}(\.{"Error:\ cannot\ creat}\)\.{e\ graphics\ context\\n}\)\.{"},\39%
\\{stderr});{}$\6
${}\\{XDestroyWindow}(\\{g\_dpy},\39\\{g\_win});{}$\6
\\{XCloseDisplay}(\\{g\_dpy});\6
${}\\{g\_dpy}\K\NULL;{}$\6
\&{return} \T{0};\6
\4${}\}{}$\2\6
${}\\{XSetForeground}(\\{g\_dpy},\39\\{g\_gc},\39\\{BlackPixel}(\\{g\_dpy},\39%
\\{DefaultScreen}(\\{g\_dpy}))){}$;\par
\U29.\fi

\M{33}The \.{WM\_DELETE\_WINDOW} protocol is the standard way for an
X~client to learn that the user has clicked the window manager's
close decoration.  We intern the atom and register it on the window;
the server will then deliver a \.{ClientMessage} carrying that atom
instead of summarily destroying the window.

\Y\B\4\X33:Register for window-manager close\X${}\E{}$\6
$\\{g\_wm\_delete}\K\\{XInternAtom}(\\{g\_dpy},\39\.{"WM\_DELETE\_WINDOW"},\39%
\\{False});{}$\6
${}\\{XSetWMProtocols}(\\{g\_dpy},\39\\{g\_win},\39{\AND}\\{g\_wm\_delete},\39%
\T{1}){}$;\par
\U29.\fi

\N{1}{34}Display teardown.

\fi

\M{35}\PB{\\{display\_close}} is the symmetric counterpart to \PB{\\{display%
\_open}}.  It
releases the graphics context, destroys the window, and closes the
display.  Each step is guarded so the function is safe to call after
a partial setup failure.

\Y\B\4\X35:Display teardown\X${}\E{}$\6
\1\1\&{static} \&{void} \\{display\_close}(\&{void})\2\2\6
${}\{{}$\1\6
\&{if} (\\{g\_dpy})\5
${}\{{}$\1\6
\&{if} (\\{g\_gc})\5
${}\{{}$\1\6
${}\\{XFreeGC}(\\{g\_dpy},\39\\{g\_gc});{}$\6
${}\\{g\_gc}\K\NULL;{}$\6
\4${}\}{}$\2\6
\&{if} (\\{g\_win})\5
${}\{{}$\1\6
${}\\{XDestroyWindow}(\\{g\_dpy},\39\\{g\_win});{}$\6
${}\\{g\_win}\K\T{0};{}$\6
\4${}\}{}$\2\6
\\{XCloseDisplay}(\\{g\_dpy});\6
${}\\{g\_dpy}\K\NULL;{}$\6
\4${}\}{}$\2\6
\4${}\}{}$\2\par
\U7.\fi

\N{1}{36}Main function.

\fi

\M{37}\PB{\\{main}} ignores its arguments --- the program takes no command-line
options.  After a successful display setup it enters a blocking event
loop driven by \PB{\\{XNextEvent}}; each event is dispatched to the
appropriate handler.  The loop terminates when either the {\bf Exit}
button is clicked or the window-manager close gesture is received.

\Y\B\4\X37:Main function\X${}\E{}$\6
\1\1\&{int} \\{main}(\&{int} \\{argc}${},\39{}$\&{char} ${}{*}{*}\\{argv})\2%
\2{}$\6
${}\{{}$\1\6
(\&{void})\,\\{argc};\6
(\&{void})\,\\{argv};\6
\&{if} ${}(\R\\{display\_open}(\,)){}$\1\5
\&{return} \T{1};\2\6
\X38:Run the event loop\X\\{display\_close}(\,);\6
\&{return} \T{0};\6
\4${}\}{}$\2\par
\U7.\fi

\M{38}The event loop dispatches three event types and silently drops the
rest.  A non-zero return from a handler breaks the loop.

\Y\B\4\X38:Run the event loop\X${}\E{}$\6
\&{for} ( ;  ; \,)\5
${}\{{}$\1\6
\\{XEvent}\\{ev};\6
${}\\{XNextEvent}(\\{g\_dpy},\39{\AND}\\{ev});{}$\7
\&{int} \\{quit}${}\K\T{0};{}$\7
\&{switch} ${}(\\{ev}.\\{type}){}$\5
${}\{{}$\1\6
\4\&{case} \\{Expose}:\6
\&{if} ${}(\\{ev}.\\{xexpose}.\\{count}\E\T{0}){}$\1\5
\\{draw\_buttons}(\,);\2\6
\&{break};\6
\4\&{case} \\{ButtonPress}:\5
${}\\{quit}\K\\{on\_button\_press}({\AND}\\{ev}.\\{xbutton});{}$\6
\&{break};\6
\4\&{case} \\{ClientMessage}:\5
${}\\{quit}\K\\{on\_client\_message}({\AND}\\{ev}.\\{xclient});{}$\6
\&{break};\6
\4\&{default}:\5
\&{break};\6
\4${}\}{}$\2\6
\&{if} (\\{quit})\1\5
\&{break};\2\6
\4${}\}{}$\2\par
\U37.\fi

\N{1}{39}Xlib API Calls.
This program uses the synchronous core of \.{libX11}.  All identifiers
listed here are declared in \.{<X11/Xlib.h>} (or \.{<X11/Xutil.h>} for
the few utility types).

\medskip
\item{$\bullet$} {\tt Display *XOpenDisplay(const char *name)}.
\par\noindent Parameter: {\tt name} --- display string such as
\.{":0"} or \.{NULL} to read \.{\$DISPLAY}.
Returns a connection handle, or \.{NULL} on failure.

\item{$\bullet$} {\tt int XCloseDisplay(Display *dpy)}.
Closes the connection and frees all server-side resources still
associated with it.

\item{$\bullet$} {\tt int DefaultScreen(Display *dpy)} (macro).
Returns the integer index of the user's default screen, used as the
second argument to many other macros.

\item{$\bullet$} {\tt Window RootWindow(Display *dpy, int scr)} (macro).
Returns the screen's root window, used as the parent of every
top-level application window.

\item{$\bullet$} {\tt unsigned long BlackPixel(Display *dpy, int scr)}
/ {\tt WhitePixel(Display *dpy, int scr)} (macros).
Return colormap entries for pure black / pure white on the named
screen.  Used here for the window border, foreground, and background.

\item{$\bullet$} {\tt Window XCreateSimpleWindow(Display *dpy,
Window parent, int x, int y, unsigned w, unsigned h, unsigned bw,
unsigned long border, unsigned long bg)}.
Allocates a new \.{InputOutput} window with the given geometry,
border width, and pixel colours.  The window is created
{\it unmapped}: it does not appear until \PB{\\{XMapWindow}}.

\item{$\bullet$} {\tt int XStoreName(Display *dpy, Window w,
const char *name)}.
Sets the \.{WM\_NAME} property; window managers display this in the
title bar.

\item{$\bullet$} {\tt int XSelectInput(Display *dpy, Window w,
long mask)}.
Specifies which event classes the server should deliver for the
window.  This program ORs three mask bits.

\item{$\bullet$} {\tt int XMapWindow(Display *dpy, Window w)}.
Requests that the window be made visible (subject to the window
manager's policy).

\item{$\bullet$} {\tt int XDestroyWindow(Display *dpy, Window w)}.
Removes the window and any descendants.

\item{$\bullet$} {\tt GC XCreateGC(Display *dpy, Drawable d,
unsigned long valuemask, XGCValues *values)}.
Allocates a graphics context.  This program passes
\PB{$\\{valuemask}\K\T{0}$} and \PB{$\\{values}\K\NULL$} to accept all server
defaults,
then customises the foreground via \PB{\\{XSetForeground}}.

\item{$\bullet$} {\tt int XSetForeground(Display *dpy, GC gc,
unsigned long pixel)}.
Sets the foreground pixel value used by subsequent draw calls.

\item{$\bullet$} {\tt int XFreeGC(Display *dpy, GC gc)}.
Releases a graphics context allocated by \PB{\\{XCreateGC}}.

\item{$\bullet$} {\tt int XDrawRectangle(Display *dpy, Drawable d,
GC gc, int x, int y, unsigned w, unsigned h)}.
Draws the {\it outline\/} of a rectangle.  The rectangle covers
$w+1$ columns and $h+1$ rows of pixels.

\item{$\bullet$} {\tt int XDrawString(Display *dpy, Drawable d,
GC gc, int x, int y, const char *string, int length)}.
Draws an 8-bit text string at the given baseline position using
the GC's current font.

\item{$\bullet$} {\tt int XQueryTextExtents(Display *dpy,
XID font\_id, const char *string, int len,
int *direction, int *ascent, int *descent,
XCharStruct *overall)}.
Asks the server for the metrics of a string in the named font.
Used here to centre the button label.

\item{$\bullet$} {\tt GContext XGContextFromGC(GC gc)} (macro).
Returns the X resource ID associated with a graphics context, used
here as the \PB{\\{font\_id}} argument to \PB{\\{XQueryTextExtents}}.

\item{$\bullet$} {\tt Atom XInternAtom(Display *dpy,
const char *name, Bool only\_if\_exists)}.
Returns the unique server-wide atom for a string property name.
Used to obtain \.{WM\_DELETE\_WINDOW}.

\item{$\bullet$} {\tt Status XSetWMProtocols(Display *dpy, Window w,
Atom *protocols, int count)}.
Tells the server which window-manager protocols this client is
willing to handle, here just \.{WM\_DELETE\_WINDOW}.

\item{$\bullet$} {\tt int XNextEvent(Display *dpy, XEvent *ev\_return)}.
Blocks until an event is available for the connection, then writes
it into the supplied union.

\fi

\M{40}{\bf POSIX System Calls and Library Functions.}
The following identifiers from the POSIX.1-2008 standard are used
directly in this program.

\medskip
\item{$\bullet$} {\tt time\_t time(time\_t *t)}.
Returns the current calendar time (seconds since the Epoch).

\item{$\bullet$} {\tt struct tm *localtime(const time\_t *clock)}.
Converts a calendar time into broken-down local time in a static
buffer.

\item{$\bullet$} {\tt size\_t strftime(char *s, size\_t max,
const char *fmt, const struct tm *tm)}.
Formats a broken-down time into a character buffer using a
\.{printf}-like format string.  This program uses
\.{"\%Y-\%m-\%d \%H:\%M:\%S"}.

\item{$\bullet$} {\tt int printf(const char *fmt, ...)}.
Writes formatted output to standard output.

\item{$\bullet$} {\tt int snprintf(char *s, size\_t n,
const char *fmt, ...)}.
Bounded \.{sprintf}; used as a fallback when \PB{\\{localtime}} fails.

\item{$\bullet$} {\tt int fputs(const char *s, FILE *stream)}.
Writes a string without a trailing newline; used for diagnostics
on \.{stderr}.

\item{$\bullet$} {\tt int fflush(FILE *stream)}.
Flushes the buffered output of \PB{\\{stream}}; called after each
{\bf Hello} click so the line appears immediately even when
\.{stdout} is piped.

\item{$\bullet$} {\tt size\_t strlen(const char *s)}.
Returns the number of bytes in the string before its null
terminator.

\item{$\bullet$} {\tt int strcmp(const char *a, const char *b)}.
Lexicographic byte-wise comparison; used to dispatch button clicks
by label.

\fi

\N{1}{41}Index.
\fi

\inx
\fin
\con
